\chapter{Structures de contrôle en PHP}

\begin{synthese}[Au programme]
Donner de l'intelligence à un script : prendre des décisions avec \code{if},
\code{else}, \code{elseif}, l'opérateur ternaire et \code{switch} ; répéter des
traitements avec \code{while}, \code{do…while}, \code{for} et \code{foreach} ;
maîtriser \code{break}, \code{continue} et la fonction \code{date()}. Ce sont les
outils qui transforment une page statique en application web dynamique.
\end{synthese}

\section{Pourquoi des structures de contrôle ?}

Jusqu'ici, nos scripts exécutaient leurs instructions \textbf{les unes après les
autres}, de haut en bas. Mais un vrai programme doit \textbf{choisir} (afficher
« Admis » ou « Recalé » selon la moyenne) et \textbf{répéter} (afficher la table de
multiplication de 7). Les \textbf{structures de contrôle} pilotent cet ordre
d'exécution.

\begin{definition}
On distingue deux grandes familles de structures de contrôle :
\begin{itemize}
  \item les structures \textbf{conditionnelles} (ou de \emph{sélection}) :
        \code{if}, \code{elseif}, \code{else}, \code{switch} ;
  \item les structures \textbf{répétitives} (ou \emph{boucles}) : \code{while},
        \code{do…while}, \code{for}, \code{foreach}.
\end{itemize}
\end{definition}

\section{Les conditions}

\subsection{La condition simple : \texttt{if}}

\begin{definition}
L'instruction \code{if} (« si ») exécute un bloc d'instructions \textbf{uniquement
si} une condition est vraie. La condition est une expression qui vaut \code{true}
(vrai) ou \code{false} (faux).
\end{definition}

\begin{syntaxe}
\code{if (condition) \{ instructions si vrai \}}
\end{syntaxe}

\begin{codephp}
<?php
  $moyenne = 12.5;
  if ($moyenne >= 10) {
    echo "Félicitations, vous êtes admis !";
  }
?>
\end{codephp}

\fig{Si la moyenne atteint 10, le message s'affiche ; sinon, rien ne se passe.}

\begin{attention}
La condition se place entre parenthèses \code{( )} et le bloc entre accolades
\code{\{ \}}. Attention à utiliser \code{==} (comparaison) et \textbf{non} \code{=}
(affectation) : \code{if (\$x = 5)} affecte 5 à \code{\$x} au lieu de tester, une
erreur très fréquente.
\end{attention}

\subsection{Choisir entre deux cas : \texttt{if…else}}

\begin{definition}
\code{else} (« sinon ») fournit un bloc à exécuter \textbf{quand la condition est
fausse}. Exactement un des deux blocs est exécuté.
\end{definition}

\begin{codephp}
<?php
  $moyenne = 8;
  if ($moyenne >= 10) {
    echo "Résultat : ADMIS";
  } else {
    echo "Résultat : RECALÉ";
  }
?>
\end{codephp}

\subsection{Choisir parmi plusieurs cas : \texttt{if…elseif…else}}

Quand il y a plus de deux possibilités, on enchaîne les conditions avec
\code{elseif}. PHP teste les conditions \textbf{dans l'ordre} et s'arrête à la
première qui est vraie.

\begin{codephp}
<?php
  // Mention au baccalauréat selon la moyenne
  $moyenne = 14.5;
  if ($moyenne >= 16) {
    echo "Mention : Très Bien";
  } elseif ($moyenne >= 14) {
    echo "Mention : Bien";
  } elseif ($moyenne >= 12) {
    echo "Mention : Assez Bien";
  } elseif ($moyenne >= 10) {
    echo "Mention : Passable";
  } else {
    echo "Ajourné";
  }
?>
\end{codephp}

\fig{Pour 14,5 : le test \texttt{>= 16} échoue, \texttt{>= 14} réussit, on affiche « Bien ».}

\begin{attention}
L'\textbf{ordre des tests compte}. Si on plaçait \code{\$moyenne >= 10} en premier,
toutes les bonnes moyennes afficheraient « Passable », car \code{14.5 >= 10} est
déjà vrai. On va donc \textbf{du plus exigeant au moins exigeant}.
\end{attention}

\subsection{L'opérateur ternaire}

\begin{definition}
L'opérateur \textbf{ternaire} \code{? :} est un raccourci pour un \code{if…else}
qui choisit \textbf{entre deux valeurs}. On le lit : « \emph{condition} ? \emph{valeur
si vrai} : \emph{valeur si faux} ».
\end{definition}

\begin{syntaxe}
\code{condition ? valeurSiVrai : valeurSiFaux}
\end{syntaxe}

\begin{codephp}
<?php
  $moyenne = 11;
  $resultat = ($moyenne >= 10) ? "Admis" : "Recalé";
  echo "Décision : $resultat";   // affiche : Décision : Admis
?>
\end{codephp}

\begin{remarque}
Le ternaire est pratique pour une affectation courte, mais pour plusieurs cas ou
des blocs longs, le \code{if…elseif…else} reste plus lisible.
\end{remarque}

\section{L'instruction \texttt{switch}}

\begin{definition}
\code{switch} compare une \textbf{même variable} à plusieurs valeurs précises
(\code{case}). C'est plus lisible qu'une longue suite de \code{elseif} lorsqu'on
teste l'égalité avec des valeurs fixes.
\end{definition}

\begin{syntaxe}
\code{switch (\$var) \{ case valeur1: ... break; case valeur2: ... break; default: ... \}}
\end{syntaxe}

\begin{codephp}
<?php
  // Afficher la filière selon un code saisi
  $code = "TI";
  switch ($code) {
    case "A4":
      echo "Série A4 (littéraire)";
      break;
    case "C":
      echo "Série C (mathématiques)";
      break;
    case "TI":
      echo "Série TI (technologies de l'information)";
      break;
    default:
      echo "Série inconnue";
  }
?>
\end{codephp}

\fig{Sans le \texttt{break}, l'exécution « tomberait » dans le case suivant.}

\begin{attention}
Le mot-clé \code{break} est \textbf{indispensable} à la fin de chaque \code{case} :
il arrête le \code{switch}. Sans lui, PHP continue d'exécuter les \code{case}
suivants (effet de « cascade »). Le bloc \code{default} est optionnel et joue le
rôle du \code{else}.
\end{attention}

\begin{exemple}
On peut volontairement \textbf{regrouper} des cas en omettant le \code{break} :
\begin{codephp}
<?php
  $jour = "samedi";
  switch ($jour) {
    case "samedi":
    case "dimanche":
      echo "C'est le week-end !";
      break;
    default:
      echo "C'est un jour de cours.";
  }
?>
\end{codephp}
\end{exemple}

\section{Opérateurs de comparaison et logiques}

Une condition repose sur des \textbf{opérateurs} qui renvoient \code{true} ou
\code{false}.

\begin{center}\small
\begin{tabular}{@{}lll@{}}
\toprule
\textbf{Opérateur} & \textbf{Signification} & \textbf{Exemple vrai}\\
\midrule
\code{==}  & égal (valeur)            & \code{5 == "5"}\\
\code{===} & identique (valeur + type)& \code{5 === 5}\\
\code{!=}  & différent                & \code{4 != 5}\\
\code{<}   & inférieur                & \code{3 < 10}\\
\code{>}   & supérieur                & \code{10 > 3}\\
\code{<=}  & inférieur ou égal        & \code{10 <= 10}\\
\code{>=}  & supérieur ou égal        & \code{12 >= 10}\\
\bottomrule
\end{tabular}
\end{center}

\begin{attention}
\code{==} compare seulement les valeurs : \code{5 == "5"} est \textbf{vrai} (PHP
convertit la chaîne en nombre). \code{===} compare \textbf{valeur ET type} :
\code{5 === "5"} est \textbf{faux} (entier contre chaîne). En cas de doute, préfère
\code{===}, plus sûr.
\end{attention}

Les \textbf{opérateurs logiques} combinent plusieurs conditions :

\begin{center}\small
\begin{tabular}{@{}lll@{}}
\toprule
\textbf{Opérateur} & \textbf{Nom} & \textbf{Vrai quand\dots}\\
\midrule
\code{\&\&} & ET  & les deux conditions sont vraies\\
\code{||}   & OU  & au moins une condition est vraie\\
\code{!}    & NON & la condition est fausse (inverse)\\
\bottomrule
\end{tabular}
\end{center}

\begin{codephp}
<?php
  $age = 17;
  $moyenne = 11;
  // Admis ET majeur ?
  if ($moyenne >= 10 && $age >= 18) {
    echo "Admis et majeur";
  }
  // Boursier si moyenne très haute OU situation sociale
  $boursier = ($moyenne >= 16) || ($age < 16);
  // Inverser une condition avec !
  if (!($moyenne >= 10)) {
    echo "Doit repasser l'examen";
  }
?>
\end{codephp}

\section{Les boucles}

\begin{definition}
Une \textbf{boucle} répète un bloc d'instructions tant qu'une condition reste vraie.
Le bloc répété s'appelle le \textbf{corps} de la boucle, et chaque répétition une
\textbf{itération}.
\end{definition}

\subsection{La boucle \texttt{while}}

\code{while} (« tant que ») teste la condition \textbf{avant} chaque itération. Si
elle est fausse dès le départ, le corps n'est jamais exécuté.

\begin{syntaxe}
\code{while (condition) \{ instructions \}}
\end{syntaxe}

\begin{codephp}
<?php
  // Compter de 1 à 5
  $i = 1;
  while ($i <= 5) {
    echo $i . " ";   // affiche : 1 2 3 4 5
    $i = $i + 1;     // ou $i++  : ne pas oublier !
  }
?>
\end{codephp}

\begin{attention}
Il faut \textbf{toujours} faire évoluer la variable de la condition (ici \code{\$i++})
à l'intérieur du corps. Sinon la condition reste vraie pour toujours : c'est une
\textbf{boucle infinie} qui bloque le serveur.
\end{attention}

\subsection{La boucle \texttt{do…while}}

\code{do…while} teste la condition \textbf{après} chaque itération : le corps est
donc exécuté \textbf{au moins une fois}, même si la condition est fausse au départ.

\begin{syntaxe}
\code{do \{ instructions \} while (condition);}
\end{syntaxe}

\begin{codephp}
<?php
  $i = 10;
  do {
    echo "Exécuté une fois, $i = $i";  // s'affiche bien
  } while ($i < 5);                     // la condition est pourtant fausse
?>
\end{codephp}

\fig{Le point-virgule après \texttt{while(...)} est obligatoire ici.}

\subsection{La boucle \texttt{for}}

\begin{definition}
La boucle \code{for} regroupe en une seule ligne les trois étapes d'une répétition
comptée : l'\textbf{initialisation}, la \textbf{condition} et l'\textbf{incrément}.
C'est la boucle idéale quand on connaît le nombre de tours à l'avance.
\end{definition}

\begin{syntaxe}
\code{for (initialisation; condition; incrément) \{ instructions \}}
\end{syntaxe}

\begin{codephp}
<?php
  // Table de multiplication de 7
  $n = 7;
  for ($i = 1; $i <= 10; $i++) {
    echo "$n x $i = " . ($n * $i) . "<br>";
  }
?>
\end{codephp}

\fig{Affiche : 7 x 1 = 7, 7 x 2 = 14, … jusqu'à 7 x 10 = 70.}

\begin{exemple}
\textbf{Somme des entiers de 1 à n} avec une boucle \code{for} :
\begin{codephp}
<?php
  $n = 100;
  $somme = 0;
  for ($i = 1; $i <= $n; $i++) {
    $somme = $somme + $i;   // ou $somme += $i;
  }
  echo "Somme de 1 à $n = $somme";   // 5050
?>
\end{codephp}
\end{exemple}

\begin{exemple}
\textbf{Compter à rebours} en faisant \textbf{décroître} le compteur :
\begin{codephp}
<?php
  for ($i = 5; $i >= 1; $i--) {
    echo $i . "... ";   // affiche : 5... 4... 3... 2... 1...
  }
  echo "Décollage !";
?>
\end{codephp}
\end{exemple}

\subsection{La boucle \texttt{foreach} (parcours de tableau)}

\begin{definition}
\code{foreach} parcourt \textbf{tous les éléments d'un tableau}, un par un, sans se
soucier des indices. C'est la boucle naturelle pour traiter une liste (élèves,
notes, produits\dots).
\end{definition}

\begin{syntaxe}
\code{foreach (\$tableau as \$valeur) \{ ... \}} \quad ou \\
\code{foreach (\$tableau as \$cle => \$valeur) \{ ... \}}
\end{syntaxe}

\begin{codephp}
<?php
  // Liste des élèves d'une classe de Tle TI
  $eleves = ["Awa", "Bello", "Chimène", "Ngono"];
  foreach ($eleves as $eleve) {
    echo "Élève : $eleve <br>";
  }
?>
\end{codephp}

Avec la forme \code{\$cle => \$valeur}, on récupère aussi l'indice (ou la clé pour
un tableau associatif) :

\begin{codephp}
<?php
  $notes = ["Maths" => 14, "PHP" => 16, "Réseaux" => 11];
  foreach ($notes as $matiere => $note) {
    echo "$matiere : $note / 20 <br>";
  }
?>
\end{codephp}

\fig{Les tableaux seront étudiés en détail au chapitre suivant.}

\section{Maîtriser les boucles : \texttt{break} et \texttt{continue}}

\begin{definition}
À l'intérieur d'une boucle :
\begin{itemize}
  \item \code{break} \textbf{sort immédiatement} de la boucle ;
  \item \code{continue} \textbf{saute le reste du tour} et passe directement à
        l'itération suivante.
\end{itemize}
\end{definition}

\begin{codephp}
<?php
  // S'arrêter dès qu'on trouve un élève recalé
  $moyennes = [12, 15, 8, 14, 16];
  foreach ($moyennes as $m) {
    if ($m < 10) {
      echo "Un élève recalé trouvé (moyenne $m), on arrête.";
      break;          // on quitte la boucle
    }
    echo "OK ($m) ";
  }
?>
\end{codephp}

\begin{codephp}
<?php
  // Afficher seulement les nombres pairs de 1 à 10
  for ($i = 1; $i <= 10; $i++) {
    if ($i % 2 != 0) {
      continue;       // impair : on saute ce tour
    }
    echo $i . " ";    // affiche : 2 4 6 8 10
  }
?>
\end{codephp}

\section{La fonction \texttt{date()}}

\begin{definition}
La fonction \code{date(\$format)} renvoie la date et l'heure \textbf{du serveur},
mises en forme selon une chaîne de \textbf{format}. Chaque lettre du format
représente un élément de la date.
\end{definition}

\begin{center}\small
\begin{tabular}{@{}lll@{}}
\toprule
\textbf{Lettre} & \textbf{Signification} & \textbf{Exemple}\\
\midrule
\code{d} & jour (2 chiffres)      & \code{26}\\
\code{m} & mois (2 chiffres)      & \code{06}\\
\code{Y} & année (4 chiffres)     & \code{2026}\\
\code{H} & heure (00 à 23)        & \code{14}\\
\code{i} & minutes (00 à 59)      & \code{05}\\
\code{N} & jour de la semaine (1=lundi) & \code{5}\\
\bottomrule
\end{tabular}
\end{center}

\begin{codephp}
<?php
  // Afficher la date et l'heure du jour
  echo "Nous sommes le " . date("d/m/Y") . "<br>";
  echo "Il est " . date("H:i") . "<br>";
?>
\end{codephp}

\begin{exemple}
\textbf{Salutation selon l'heure.} On combine \code{date("H")} (l'heure sous forme
de nombre) et un \code{if…elseif…else} :
\begin{codephp}
<?php
  $heure = (int) date("H");   // ex. 14 ; (int) force un entier
  if ($heure < 12) {
    echo "Bonjour et bonne matinée !";
  } elseif ($heure < 18) {
    echo "Bon après-midi !";
  } else {
    echo "Bonne soirée !";
  }
?>
\end{codephp}
\end{exemple}

\begin{exemple}
\textbf{Afficher le nom du jour} grâce à \code{date("N")} et un tableau :
\begin{codephp}
<?php
  $jours = [1 => "Lundi", "Mardi", "Mercredi", "Jeudi",
            "Vendredi", "Samedi", "Dimanche"];
  $numero = (int) date("N");          // 1 = lundi … 7 = dimanche
  echo "Aujourd'hui, c'est " . $jours[$numero];
?>
\end{codephp}
\end{exemple}

\begin{methode}
\textbf{Construire une condition correcte.} (1)~Identifier la donnée à tester
(moyenne, âge, quantité\dots) ; (2)~choisir le bon opérateur de comparaison ;
(3)~ordonner les cas du plus précis au plus général ; (4)~vérifier qu'\textbf{un
seul} bloc s'exécute pour chaque valeur d'entrée.
\end{methode}

\section{Exercices}

\rubrique{Application}

\exo{Pair ou impair}\diff{1}
Une variable \code{\$nombre} contient un entier. Écris un script qui affiche
« pair » si le nombre est pair, « impair » sinon. (Indice : un nombre est pair si
\code{\$nombre \% 2 == 0}.)

\exo{Admission}\diff{1}
La variable \code{\$moyenne} contient la moyenne d'un élève. Affiche « Admis » si
elle vaut au moins 10, sinon « Recalé », en utilisant l'\textbf{opérateur ternaire}.

\exo{Le plus grand de trois nombres}\diff{2}
Trois notes sont rangées dans \code{\$a}, \code{\$b} et \code{\$c}. Affiche la plus
grande des trois à l'aide de conditions.

\exo{Compte à rebours}\diff{1}
À l'aide d'une boucle \code{for}, affiche un compte à rebours de 10 jusqu'à 1, puis
le mot « Partez ! ».

\exo{Table de multiplication}\diff{2}
La variable \code{\$n} contient un entier. Affiche sa table de multiplication
(de \code{\$n}~$\times$~1 à \code{\$n}~$\times$~10) à l'aide d'une boucle.

\rubrique{Entraînement}

\exo{Prix avec remise selon la quantité}\diff{2}
Dans une librairie de Yaoundé, un cahier coûte 500 FCFA. Si le client achète 10
cahiers ou plus, il bénéficie de 10\,\% de remise ; à partir de 50 cahiers, la
remise est de 20\,\%. La quantité est dans \code{\$qte}. Calcule et affiche le prix
total à payer en francs CFA.

\exo{Mention au baccalauréat}\diff{2}
La variable \code{\$moyenne} contient une moyenne sur 20. Affiche la mention :
« Très Bien » ($\geq 16$), « Bien » ($\geq 14$), « Assez Bien » ($\geq 12$),
« Passable » ($\geq 10$) ou « Ajourné » sinon.

\exo{Table de multiplication en HTML}\diff{2}
Reprends la table de multiplication de \code{\$n}, mais affiche-la dans un véritable
\textbf{tableau HTML} (balises \code{<table>}, \code{<tr>}, \code{<td>}).

\exo{Somme et moyenne de n notes}\diff{3}
On dispose d'un tableau \code{\$notes} contenant les notes d'un élève (par exemple
\code{[12, 15, 8, 14, 16]}). Calcule et affiche la \textbf{somme} et la
\textbf{moyenne} de ces notes, ainsi que le verdict « Admis » ou « Recalé ».

\exo{FizzBuzz scolaire}\diff{3}
Pour les nombres de 1 à 30, affiche : « Récré » si le nombre est divisible par 3,
« Pause » s'il est divisible par 5, « Récré-Pause » s'il est divisible par 3 ET
par 5, et le nombre lui-même sinon.

\exo{Année bissextile}\diff{3}
La variable \code{\$annee} contient une année. Affiche si elle est bissextile.
(Règle : une année est bissextile si elle est divisible par 4 \textbf{et} non
divisible par 100, \textbf{ou} si elle est divisible par 400.)

\section{Corrigés}

\corrige{de l'exercice 1}
\begin{codephp}
<?php
  $nombre = 7;
  if ($nombre % 2 == 0) {
    echo "$nombre est pair";
  } else {
    echo "$nombre est impair";
  }
?>
\end{codephp}

\corrige{de l'exercice 2}
\begin{codephp}
<?php
  $moyenne = 11.5;
  $resultat = ($moyenne >= 10) ? "Admis" : "Recalé";
  echo $resultat;
?>
\end{codephp}

\corrige{de l'exercice 3}
\begin{codephp}
<?php
  $a = 12; $b = 17; $c = 9;
  $max = $a;                 // on suppose que $a est le plus grand
  if ($b > $max) {
    $max = $b;
  }
  if ($c > $max) {
    $max = $c;
  }
  echo "La plus grande note est : $max";
?>
\end{codephp}

\corrige{de l'exercice 4}
\begin{codephp}
<?php
  for ($i = 10; $i >= 1; $i--) {
    echo $i . " ";
  }
  echo "Partez !";
?>
\end{codephp}

\corrige{de l'exercice 5}
\begin{codephp}
<?php
  $n = 8;
  for ($i = 1; $i <= 10; $i++) {
    echo "$n x $i = " . ($n * $i) . "<br>";
  }
?>
\end{codephp}

\corrige{de l'exercice 6}
\begin{codephp}
<?php
  $prixUnitaire = 500;   // FCFA
  $qte = 60;
  if ($qte >= 50) {
    $remise = 0.20;
  } elseif ($qte >= 10) {
    $remise = 0.10;
  } else {
    $remise = 0;
  }
  $total = $prixUnitaire * $qte * (1 - $remise);
  echo "Total à payer : " . $total . " FCFA";
?>
\end{codephp}

\corrige{de l'exercice 7}
\begin{codephp}
<?php
  $moyenne = 14.5;
  if ($moyenne >= 16) {
    echo "Très Bien";
  } elseif ($moyenne >= 14) {
    echo "Bien";
  } elseif ($moyenne >= 12) {
    echo "Assez Bien";
  } elseif ($moyenne >= 10) {
    echo "Passable";
  } else {
    echo "Ajourné";
  }
?>
\end{codephp}

\corrige{de l'exercice 8}
\begin{codephp}
<?php
  $n = 6;
  echo "<table border='1'>";
  for ($i = 1; $i <= 10; $i++) {
    echo "<tr>";
    echo "<td>$n x $i</td>";
    echo "<td>" . ($n * $i) . "</td>";
    echo "</tr>";
  }
  echo "</table>";
?>
\end{codephp}

\corrige{de l'exercice 9}
\begin{codephp}
<?php
  $notes = [12, 15, 8, 14, 16];
  $somme = 0;
  foreach ($notes as $note) {
    $somme += $note;
  }
  $moyenne = $somme / count($notes);   // count() compte les éléments
  echo "Somme : $somme <br>";
  echo "Moyenne : " . round($moyenne, 2) . " / 20 <br>";
  echo ($moyenne >= 10) ? "Admis" : "Recalé";
?>
\end{codephp}

\fig{La fonction \texttt{count()} renvoie le nombre d'éléments d'un tableau.}

\corrige{de l'exercice 10}
\begin{codephp}
<?php
  for ($i = 1; $i <= 30; $i++) {
    if ($i % 3 == 0 && $i % 5 == 0) {
      echo "Récré-Pause <br>";
    } elseif ($i % 3 == 0) {
      echo "Récré <br>";
    } elseif ($i % 5 == 0) {
      echo "Pause <br>";
    } else {
      echo $i . " <br>";
    }
  }
?>
\end{codephp}

\fig{On teste d'abord le cas « divisible par 3 ET 5 », sinon il ne serait jamais atteint.}

\corrige{de l'exercice 11}
\begin{codephp}
<?php
  $annee = 2024;
  if (($annee % 4 == 0 && $annee % 100 != 0) || ($annee % 400 == 0)) {
    echo "$annee est bissextile";
  } else {
    echo "$annee n'est pas bissextile";
  }
?>
\end{codephp}

\begin{aretenir}
Un script intelligent \textbf{décide} (\code{if}, \code{elseif}, \code{else},
ternaire, \code{switch}) et \textbf{répète} (\code{while}, \code{do…while},
\code{for}, \code{foreach}). Ordonne tes conditions du plus précis au plus général,
n'oublie jamais le \code{break} dans un \code{switch} ni de faire évoluer le
compteur d'une boucle, et choisis \code{for} quand tu connais le nombre de tours,
\code{while} sinon.
\end{aretenir}
